\chapter{Material suplementario}
\label{app:supplementary}

\section{Enlaces del proyecto}

\begin{table}[H]
  \centering
  \caption{Recursos del proyecto Litoteca}
  \label{tab:recursos}
  \begin{tabular}{@{}p{4.2cm}p{8.2cm}@{}}
    \toprule
    Recurso & URL / referencia \\
    \midrule
    Repositorio GitHub & \url{https://github.com/JheisonMB/litoteca} \\
    Aplicación desplegada & \url{https://litoteca.pages.dev/} \\
    Mock de fichas & \url{https://litoteca.pages.dev/?mock=1} \\
    Galería de fichas & \url{https://litoteca.pages.dev/?galeria=1} \\
    Dataset fuente & \url{https://www.kaggle.com/datasets/stealthtechnologies/rock-classification} \\
    Custom Vision & \url{https://azure.microsoft.com/services/cognitive-services/custom-vision-service/} \\
    \bottomrule
  \end{tabular}
\end{table}

\section{Trazabilidad del dataset}

El dataset preparado incluye los siguientes artefactos en \texttt{data/prepared/}:

\begin{itemize}
  \item \texttt{manifest.csv}: procedencia, split, clase, hash perceptual.
  \item \texttt{stats.json}: conteos por clase y motivos de descarte.
  \item \texttt{conflicts.csv}: 22 pares con etiqueta contradictoria.
  \item \texttt{near\_duplicates.csv}: casi-duplicados entre clusters.
  \item \texttt{descartadas.csv}: registros excluidos.
\end{itemize}

Para el entrenamiento limpio se generó \texttt{data/prepared-clean-train/} con una
imagen por cluster y exclusión de conflictos de etiqueta. El reporte de limpieza se
encuentra en \texttt{clean\_report.json}.

\section{Evidencia de despliegue}

El despliegue se realiza con GitHub Actions (\texttt{.github/workflows/deploy.yml})
y Cloudflare Wrangler. Los secretos del modelo se configuran en Cloudflare Pages
como variables de runtime de la Function, no en el frontend:

\begin{itemize}
  \item \texttt{CUSTOM\_VISION\_ENDPOINT}
  \item \texttt{CUSTOM\_VISION\_KEY}
  \item \texttt{CUSTOM\_VISION\_PROJECT}
  \item \texttt{CUSTOM\_VISION\_ITERATION}
\end{itemize}
